\documentclass{article}

\usepackage{amsmath} % math stuff
\usepackage{amssymb} % math stuff
\usepackage{array} % equations and stuff
\usepackage{bm} % bold math
%\usepackage{booktabs} % extra table rule options
%\usepackage{caption} % suppressed table numbering; incompatible with revtex, and longtable, I think
\usepackage{comment} % comment environment
%\usepackage{enumitem} % customization of enumeration, itemize, and description
\usepackage[T1]{fontenc} % font encoding for special characters, must also use scalable font package
\usepackage[margin=0.8in]{geometry} % paper sizes and margins (but be careful not to mess up pre-defined pages)
\usepackage{graphicx} % for graphics
%\usepackage{helvet} % default font is the helvetica postscript font
\usepackage[utf8]{inputenc} % special characters in tex input
\usepackage{layouts} % print units like widths
\usepackage{lipsum} % lorem ipsum filler text
\usepackage{lmodern} % scalable font?
\usepackage{longtable} % multi-page tables
\usepackage{makecell} % specify line-breaks in table cells
\usepackage{mathrsfs} % math script font
\usepackage{mhchem} % easier chemical formula
\usepackage{microtype} % allows disabling of ligatures
\usepackage{multicol} % multicolumns
%\usepackage{newcent} % new century schoolbook font
\usepackage{nicefrac}
\usepackage{numprint} % print and format (large) numbers
\usepackage{parskip} % removes paragraph indentation, and adjusts paragraph skip, as well as list items
\usepackage{pdfpages} % add pdf files as pages
%\usepackage{setspace} % adjust text spacing and indents
\usepackage{siunitx} % decimal alignment, and degree symbol
\usepackage{subfigure} % divided figures
%\usepackage{tabu} % extra table options
\usepackage{textcomp} % symbols
\usepackage{threeparttablex} % better footnotes with longtable
\usepackage{titling} % title placement
\usepackage{ulem} % strikethrough text
\usepackage{upgreek} % upright Greek letters
%\usepackage{url} % superceded by hyperref
\usepackage{verbatim} % verbatim environment
\usepackage{xcolor} % colors and color boxes
\usepackage{xspace} % commands that don't eat up white space
\usepackage{hyperref} % links and page setup; should always come last

\hypersetup{
 bookmarks=true,
 colorlinks=true,
 citecolor=blue,
 linkcolor=blue,
 urlcolor=blue,
 pdfstartview={XYZ null null 1.0} % default open view is 100%
}

\DisableLigatures[f,t]{encoding = T1} % disable ff, fi, fl, tt ligatures; without options, it also disables -- = endash
\renewcommand{\arraystretch}{1.0} % extra vertical (and horizontal?) space in tables

% define centered, left- and right-aligned columns with specified widths
\newcommand{\PreserveBackslash}[1]{\let\temp=\\#1\let\\=\temp}
\newcolumntype{C}[1]{>{\PreserveBackslash\centering}p{#1}}
\newcolumntype{L}[1]{>{\PreserveBackslash\raggedright}p{#1}}
\newcolumntype{R}[1]{>{\PreserveBackslash\raggedleft}p{#1}}

\begin{document}

\pagestyle{empty} % don't number pages

% custom title
\begin{center}
{\LARGE Express Riddler}

\vspace{0.15in}

{\Large 18 March 2022}
\end{center}


\section*{Riddle:}

Earlier today, James's boss was surprised to find out that not only did no one on their team have a birthday this week, but that nobody was celebrating a birthday for the entire \textit{month}.
With a total of 40 people on the team, the probability of this happening seemed to be miniscule.

But was that really the case?
What was the probability that none of the 40 people had birthdays this month?
(For the purpose of this riddle, assume that a year consists of 12 equally long months.
It's a sufficiently good approximation!)

\textit{Extra credit}: What is the probability that there is at least one month in the year during which none of the 40 people had birthdays (not necessarily \textit{this} month)?


\section*{Solution:}

The probability that a single person doesn't have a birthday in March (or any other single month) is $\nicefrac{11}{12}\approx0.833$.
For 40 people with independent and identically distributed birthdays, the combined probability is $(\nicefrac{11}{12})^{12}\approx0.0308$.
So the solution to the first question is (approximately) \fcolorbox{red}{white}{\textbf{3.08\%}}\,.

The probability $P$ that there is at least one month with no birthdays is

\begin{align*}
P &= \sum_{n=1}^{12}(-1)^{n+1}\binom{12}{n}\left(\frac{n}{12}\right)^{40} \\
  &= \binom{12}{1}\left(\frac{1}{12}\right)^{40} - \binom{12}{2}\left(\frac{2}{12}\right)^{40} + \binom{12}{3}\left(\frac{3}{12}\right)^{40} - \binom{12}{4}\left(\frac{4}{12}\right)^{40} +\dots
\end{align*}

As discussed above, the first term in the sum ($(\nicefrac{11}{12})^{40}$) is the probability that there are no birthdays in any particular month.
This must be added 12 times (12 choose 1) for each month.
Adding the first two terms actually overcounts situations in which there are two different months with no birthdays.
So to correct for that, the second term in the sum ($(\nicefrac{10}{12})^{40}$) is the probability that there are no birthdays in any two particular months.
There are 66 (12 choose 2) pairs of months that must be included in this correction.
Similarly, this overcorrects (and undercounts) situations in which there are three different months with no birthdays.
So the third term in the sum ($(\nicefrac{9}{12})^{40}$) is the probability that there are no birthdays in any two particular months.
And so on, and so on.
The twelfth term in the sum ($(\nicefrac{0}{12})^{40}$) is the probability that there are no birthdays in all 12 months.
This is of course 0, since each birthday must occur in some month.

Adding all of these terms gives the (approximate) probability of \fcolorbox{red}{white}{\textbf{32.68\%}}\,.


\end{document}